%% ========================================================
\section{Електронні переходи та УФ-видимі спектри}\label{sec:electronic}
%% ========================================================

%% --------------------------------------------------------
\subsection{Теорія збуджених станів}
%% --------------------------------------------------------

Поглинання фотона молекулою призводить до переходу електрона
з основного електронного стану $S_0$ у збуджений $S_n$.
Енергія переходу визначається різницею повних енергій:
\[
    \hbar\omega_n = E_n - E_0.
\]

Однак у реальному спектрі спостерігаються не «чисті» електронні переходи,
а \emph{електронно-вібраційні}, тобто такі, де відбувається зміна не лише електронного,
а й вібраційного стану молекули:
\[
    |S_0, v''\rangle \;\longrightarrow\; |S_n, v'\rangle.
\]

Через різницю форм потенціальних кривих $E_0(R)$ і $E_n(R)$
мінімальні відстані рівноваги $R_e^{(0)}$ і $R_e^{(n)}$ не збігаються.
У момент поглинання фотона ядра практично не встигають зміститись —
це \textbf{принцип Франка–Кондона}:
електронний перехід відбувається при фіксованих координатах ядер.
На діаграмі потенціальних кривих це означає \emph{вертикальний перехід}:
\[
    R_{\text{збудж.}} = R_{\text{основ.}}.
\]

%---------------------------------------------------------
\begin{figure}[h!]\centering
    %\includegraphics[width=0.8\linewidth]{\currfiledir/franck_condon_diagram.pdf}
    \caption{Схематичне зображення принципу Франка–Кондона.
    Вертикальні переходи $S_0 \to S_n$ відбуваються без зміни координат ядер.}
    \label{pic:franck_condon}
\end{figure}
%---------------------------------------------------------

Згідно з наближенням Борна–Оппенгаймера,
повна хвильова функція молекули розкладається на електронну та ядерну частини:
\[
\Psi_{n,v}(r,R) = \psi_n(r;R)\,\chi_{n,v}(R),
\]
де $\psi_n(r;R)$ описує електрони, а $\chi_{n,v}(R)$ — ядерний (вібраційний) рух.

Імовірність переходу між станами $|S_0,v''\rangle$ і $|S_n,v'\rangle$
визначається квадратом модуля дипольного матричного елемента:
\[
I_{v'\!\leftarrow\!v''} \propto
\left| \langle \Psi_{n,v'} | \hat{\boldsymbol{\mu}} | \Psi_{0,v''} \rangle \right|^2.
\]

Якщо електронний оператор $\hat{\boldsymbol{\mu}}$ слабко залежить від координат ядер,
вираз розкладається на електронну та вібраційну частини:
\[
I_{v'\!\leftarrow\!v''} \propto
\underbrace{
\left| \langle \psi_n | \hat{\boldsymbol{\mu}}_e | \psi_0 \rangle \right|^2
}_{\text{електронний перехід}}
\underbrace{
\left| \langle \chi_{n,v'} | \chi_{0,v''} \rangle \right|^2
}_{\text{фактор Франка–Кондона }F_{v'\!,v''}}.
\]

Другий множник
\[
F_{v'\!,v''} = \left| \langle \chi_{n,v'} | \chi_{0,v''} \rangle \right|^2
\]
називається \textbf{фактором Франка–Кондона}
і відображає перекриття хвильових функцій вібраційних рівнів.
Він визначає, які з переходів $v''\!\to v'$ будуть найінтенсивнішими:
чим сильніше перекриття, тим більша інтенсивність.

---

\subsubsection*{Зв’язок із силою осцилятора}

\textbf{Сила осцилятора} є мірою інтенсивності переходу в спектрі поглинання.
Для кожного переходу $|S_0,v''\rangle \to |S_n,v'\rangle$ вона записується як:
\[
f_{v'\!,v''} = \frac{2}{3}\,\omega_{v'\!,v''}
\left| \langle \Psi_{n,v'} | \hat{\boldsymbol{\mu}} | \Psi_{0,v''} \rangle \right|^2.
\]
Підставляючи розклад із принципу Франка–Кондона, отримаємо:
\[
f_{v'\!,v''} \propto
\underbrace{\left| \langle \psi_n | \hat{\boldsymbol{\mu}}_e | \psi_0 \rangle \right|^2}_{\text{тип переходу}}
\times
\underbrace{F_{v'\!,v''}}_{\text{вібраційне перекриття}}.
\]

Отже:
\begin{itemize}
\item електронна частина визначає \emph{загальну інтенсивність переходу};
\item вібраційна частина (фактор Франка–Кондона) \emph{розподіляє цю інтенсивність між коливальними рівнями}.
\end{itemize}

Якщо потенціальні криві $E_0(R)$ і $E_n(R)$ мають близькі мінімуми,
то головним буде перехід $v''=0 \to v'=0$,
і спектр має вузьку, симетричну смугу.
При великому зсуві $R_e^{(n)}-R_e^{(0)}$ інтенсивність розподіляється
по багатьох вібраційних лініях, утворюючи типову коливальну структуру у смузі поглинання.


%% --------------------------------------------------------
\subsection{TDDFT розрахунок для формальдегіду}
%% --------------------------------------------------------

Розглянемо молекулу формальдегіду \ce{H2CO} як приклад:

\inputcode{h2co_tddft.py}

\textbf{Аналіз збуджень для \ce{H2CO}:}
\begin{center}
    \begin{tabular}{lcccc}
        \toprule
        Перехід & Тип           & $\lambda$ (нм) & $f$   & Характер       \\
        \midrule
        $S_1$   & $n\to\pi^*$   & 355            & 0.001 & Заборонений    \\
        $S_2$   & $\pi\to\pi^*$ & 185            & 0.152 & Дозволений     \\
        $S_3$   & $n\to 3s$     & 172            & 0.021 & Рідбергівський \\
        \bottomrule
    \end{tabular}
\end{center}

\textbf{Фізична інтерпретація:}
\begin{itemize}
    \item $n\to\pi^*$: неподілена пара O $\to$ антизв'язуюча МО C=O
    \item Низька сила осцилятора через симетрію
    \item $\pi\to\pi^*$: основна смуга поглинання в УФ
    \item Енергії залежать від функціоналу DFT
\end{itemize}

%% --------------------------------------------------------
\subsection{Вибір функціоналу для TDDFT}
%% --------------------------------------------------------

Різні функціонали дають різну точність для збуджень:

\inputcode{formaldehyde_functional_comparison.py}

\textbf{Порівняння функціоналів (перехід $n\to\pi^*$):}
\begin{center}
    \begin{tabular}{lcc}
        \toprule
        Функціонал     & $\lambda$ (нм) & Похибка (нм) \\
        \midrule
        B3LYP          & 355            & +25          \\
        PBE0           & 342            & +12          \\
        CAM-B3LYP      & 335            & +5           \\
        $\omega$B97X-D & 332            & +2           \\
        \midrule
        Експеримент    & 330            & ---          \\
        \bottomrule
    \end{tabular}
\end{center}

\textbf{Рекомендації:}
\begin{itemize}
    \item Гібридні функціонали краще за чисті GGA
    \item Range-separated (CAM-B3LYP, $\omega$B97X-D) найточніші
    \item Для переносу заряду обов'язково range-separated
    \item B3LYP часто завищує довжини хвиль
\end{itemize}

%% --------------------------------------------------------
\subsection{Візуалізація спектру}
%% --------------------------------------------------------

Для побудови спектру використовуємо гаусові або лоренцеві контури:
\[
    \varepsilon(\lambda) = \sum_n f_n \cdot \frac{1}{\sigma\sqrt{2\pi}}
    \exp\left[-\frac{(\lambda - \lambda_n)^2}{2\sigma^2}\right]
\]

\inputcode{plot_uv_spectrum.py}

%---------------------------------------------------------
\begin{figure}[h!]\centering
    %\includegraphics[width=\linewidth]{\currfiledir/h2co_uv_spectrum.pdf}
    \caption{УФ-спектр формальдегіду, розрахований методом TDDFT/CAM-B3LYP/aug-cc-pVDZ.}
    \label{pic:h2co_uv_spectrum}
\end{figure}
%---------------------------------------------------------

\textbf{Параметри уширення:}
\begin{itemize}
    \item Типове $\sigma = 0.3$--$0.5$ eV для конденсованої фази
    \item $\sigma = 0.1$--$0.2$ eV для газової фази
    \item Експериментальне уширення враховує розподіл за $T$
\end{itemize}

%% --------------------------------------------------------
\subsection{Аналіз характеру переходів}
%% --------------------------------------------------------

TDDFT надає інформацію про орбіталі, задіяні у переході:

\inputcode{analyze_transitions.py}

\textbf{Приклад виводу для $S_1$ стану \ce{H2CO}:}
\begin{minted}{text}
Excited State 1: 3.492 eV (355 nm)  f=0.0012
    HOMO-1 -> LUMO     0.11 (1.2%)
    HOMO   -> LUMO     0.69 (47.6%)
    HOMO   -> LUMO+1   0.08 (0.6%)
\end{minted}

\textbf{Інтерпретація:}
\begin{itemize}
    \item Основний внесок: HOMO $\to$ LUMO (48\%)
    \item HOMO --- неподілена пара n(O)
    \item LUMO --- антизв'язуюча $\pi^*$(C=O)
    \item Тип переходу: $n\to\pi^*$
\end{itemize}